\begin{ouverture}[lunette]{18}{Lentilles et vision}{Ondes et signaux}{Observer · Modéliser · Expérimenter · Expliquer}
\objectif{Identifier centre optique, foyers et distance focale.}
\objectif{Construire une image réelle avec des rayons remarquables.}
\objectif{Déterminer un grandissement et modéliser l’œil.}
\prerequis{Proportionnalité, conversions et lecture de graphiques.}
\end{ouverture}
\essentiel
\essSection{1}{Une lentille mince convergente}
Une lentille convergente dévie un faisceau parallèle à l'axe optique vers son foyer image $F^{\prime}$. Son centre optique est O ; les foyers F et $F^{\prime}$ sont symétriques par rapport à O. La distance focale $f'$ est la distance positive $OF^{\prime}$.
\schema[corrige]{lentille-foyers}
\begin{propriete}[Trois rayons utiles]Un rayon passant par O n'est pas dévié dans le modèle de la lentille mince. Un rayon incident parallèle à l'axe ressort en passant par $F^{\prime}$. Un rayon passant par F ressort parallèle à l'axe.\end{propriete}
\essSection{2}{Construire une image réelle}
Pour un objet AB perpendiculaire à l'axe, avec A sur l'axe et placé avant le foyer objet, tracer au moins deux rayons issus de B. Leur intersection après la lentille donne $B^{\prime}$ ; $A^{\prime}$ est le pied sur l'axe. Tous les rayons issus de B et transmis dans le modèle convergent en $B^{\prime}$, pas seulement les deux rayons tracés.
\schema[corrige]{lentille-2f}
Une image réelle peut être recueillie sur un écran placé dans son plan. Elle est ici renversée. Si l'objet est trop près de la lentille, la construction d'une image réelle sur écran n'est plus possible : le cas de la loupe constitue un autre contexte.
\essSection{3}{Grandissement et œil réduit}
Le grandissement algébrique est $\gamma=\overline{A^{\prime}B^{\prime}}/\overline{AB}$ : il est négatif pour l'image renversée. Sa valeur absolue compare les tailles. Par Thalès, pour les longueurs positives objet-lentille $p$ et lentille-image $p^{\prime}$, $|\gamma|=p^{\prime}/p$ dans cette construction.
\schema[corrige]{oeil}
L'œil réduit associe un diaphragme à la pupille, une lentille convergente à l'ensemble optique cornée-cristallin, et un écran à la rétine. L'image nette se forme sur la rétine. L'accommodation modifie la convergence du système optique ; le cerveau interprète l'image reçue.
\essSection{4}{Expérimenter et modéliser}
\begin{experience}[Former une image]Aligner un objet lumineux adapté, une lentille et un écran sur un banc. Repérer la distance focale avec un objet lointain, puis placer l'objet au-delà de F. Déplacer l'écran jusqu'à la netteté ; mesurer tailles et distances, comparer la construction et le grandissement. Ne jamais viser le Soleil.\end{experience}
\begin{savoirfaire}[Un dessin exploitable]Indiquer l'échelle spatiale, O, F, $F^{\prime}$, AB, l'axe et le sens de propagation. Les flèches des rayons indiquent la propagation ; les flèches de AB et $A^{\prime}B^{\prime}$ indiquent l'orientation de l'objet et de l'image.\end{savoirfaire}
\exercices
\rubrique{Prendre confiance}
\begin{blocLivre}{ex-18-1}\begin{exercice}[1]{Repérer les foyers}
La lentille a une distance focale de 10 cm.
\schema{lentille-foyers}
\begin{questions}
\item À quelle distance de O placer F et $F^{\prime}$ ?
\item Quel foyer reçoit les rayons arrivant parallèlement à l’axe ?
\end{questions}
\end{exercice}\end{blocLivre}
\begin{blocLivre}{ex-18-2}\begin{exercice}[1]{Le rayon central}
Un rayon passe par le centre optique O.
\begin{questions}
\item Comment le tracer après la lentille dans le modèle mince ?
\item Ce rayon suffit-il à localiser seul l’image d’un point ?
\end{questions}
\end{exercice}\end{blocLivre}
\begin{blocLivre}{ex-18-3}\begin{exercice}[1]{Un rayon parallèle}
Un rayon issu de B atteint la lentille parallèlement à l’axe.
\begin{questions}
\item Par quel point passe le rayon émergent ?
\item Quel rayon émergent serait parallèle à l’axe ?
\end{questions}
\end{exercice}\end{blocLivre}
\begin{blocLivre}{ex-18-4}\begin{exercice}[1]{Lire une image réelle}
Un objet de 2,0 cm placé à 20 cm d’une lentille de focale 10 cm est représenté à l’échelle indiquée.
\schema{lentille-2f}
\begin{questions}
\item Construire l’image avec deux rayons.
\item Donner sa position, sa taille et son sens.
\end{questions}
\end{exercice}\end{blocLivre}
\rubrique{S’entraîner}
\begin{blocLivre}{ex-18-5}\begin{exercice}[2]{Un grandissement}
Un objet de 3,0 cm donne une image renversée de 1,5 cm.
\begin{questions}
\item Calculer le grandissement algébrique.
\item L’image est-elle agrandie ou réduite ?
\end{questions}
\end{exercice}\end{blocLivre}
\begin{blocLivre}{ex-18-6}\begin{exercice}[2]{Construire une image réduite}
La focale vaut 10 cm, l’objet est à 30 cm de O et mesure 2,0 cm. La figure est à l’échelle.
\schema{lentille-3f}
\begin{questions}
\item Construire et lire la position de l’image.
\item Lire sa taille et calculer $\gamma$.
\end{questions}
\end{exercice}\end{blocLivre}
\begin{blocLivre}{ex-18-7}\begin{exercice}[2]{Le modèle de l’œil}
Utiliser le schéma de l’œil réduit.
\schema{oeil}
\begin{questions}
\item Associer pupille, cristallin et rétine aux trois éléments du modèle.
\item Où doit se former l’image pour une vision nette ?
\end{questions}
\end{exercice}\end{blocLivre}
\begin{blocLivre}{ex-18-8}\begin{exercice}[2]{Utiliser Thalès}
Une image réelle est à 24 cm d’une lentille ; l’objet est à 12 cm et mesure 1,0 cm.
\begin{questions}
\item Déterminer $|\gamma|$.
\item Donner la taille et le sens de l’image.
\end{questions}
\end{exercice}\end{blocLivre}
\rubrique{Raisonner et approfondir}
\begin{blocLivre}{ex-18-9}\begin{exercice}[3]{Un écran mal placé}
Une image nette devrait se former à 15 cm après la lentille ; l’écran est à 20 cm.
\begin{questions}
\item Pourquoi l’image observée est-elle floue ?
\item Doit-on changer la position de l’écran pour ce montage fixe ?
\end{questions}
\end{exercice}\end{blocLivre}
\begin{blocLivre}{ex-18-10}\begin{exercice}[3]{Cacher une moitié}
On masque une moitié de la lentille, sans déplacer le montage.
\begin{questions}
\item Une moitié de l’image disparaît-elle nécessairement ?
\item Quel effet principal peut-on attendre sur la luminosité ?
\end{questions}
\end{exercice}\end{blocLivre}
\begin{blocLivre}{ex-18-11}\begin{exercice}[3]{Une accommodation}
Un objet se rapproche de l’œil alors que la distance entre le système optique et la rétine reste pratiquement fixe.
\begin{questions}
\item Quelle propriété l’œil peut-il modifier pour garder une image nette ?
\item Le modèle demande-t-il de déplacer la rétine comme un écran de banc ?
\end{questions}
\end{exercice}\end{blocLivre}
\begin{blocLivre}{ex-18-12}\begin{exercice}[3]{Valider un tracé}
Un dessin place l’objet à 30 cm, l’image à 15 cm, et donne $|\gamma|=2$.
\begin{questions}
\item Détecter l’incohérence avec Thalès.
\item Quelle valeur doit-on obtenir pour une image réelle renversée ?
\end{questions}
\end{exercice}\end{blocLivre}

